% ============================================================================
%  Capítulo 7 — Discusión
% ============================================================================
\chapter{Discusión}
\label{cap:discusion}

\section{Interpretación de los hallazgos}
\label{sec:hipotesis_discusion}

\paragraph{Dinámica de la planta.} Al variar el $b_0$ del controlador entre $1/6$ y $6$ veces el
nominal, se evidencia que el lazo de control es capaz de soportar las discrepancias entre las
ganancias (Figuras~\ref{fig:adrc_variacion}(d) y \ref{fig:sat_adrc_gpi_variacion}(d)) sin entrar en
inestabilidad. Por lo que, incluso sin una ganancia exacta, el sistema de control tiene la suficiente
robustez para seguir funcionando, lo que es importante en esta planta, donde la ganancia cambia de una
ventana a otra (Sección~\ref{sec:ident_arx}). Ahora bien, esto es sin retardo, donde el margen de
ganancia es infinito (Figura~\ref{fig:adrc_margenes_retardo}). Con el retardo de diseño el margen de
ganancia es de $2{,}12$ (Figura~\ref{fig:adrc_margenes_retardo}), así que el lazo se vuelve inestable
si el $b_0$ del controlador queda por debajo de $b_{0,\mathrm{planta}}/2{,}12$. Por eso, si hay duda, conviene que $b_0$ quede por encima del de
la planta: el lazo pierde velocidad, pero no robustez. La elección de $b_0$ es uno de los aspectos
prácticos del ADRC que ya sistematizan Herbst \cite{herbst2013practical} y Madoński y Herman
\cite{madonski2015survey}; aquí se acota con el rango que entrega la identificación de la planta.

\paragraph{Retardo.} El sistema es capaz de aguantar las variaciones del retardo de la planta,
incluso hasta el doble del de diseño, sin desestabilizarse (Figura~\ref{fig:sat_adrc_gpi_variacion}(c)).
La cuestión es que se ve un incremento en el sobrepico y mayores oscilaciones ante una referencia tipo
paso, porque con el retardo el margen de fase baja (Figura~\ref{fig:adrc_margenes_retardo}), y el
margen de fase se relaciona con el amortiguamiento del lazo: entre menos margen, más sobrepico. Con
Bézier y con el prefiltro el sobrepico casi desaparece (Figuras~\ref{fig:adrc_bezier_retardo} y
\ref{fig:pref_escalon}). El margen de fase del lazo no cambia, lo que cambia es que la referencia ya no
excita esa oscilación. Herbst y Madoński proponen suavizar la referencia del ADRC basado en el error
para evitar los esfuerzos de control muy grandes ante un escalón \cite[fig.~6.13 y
sec.~6.4.2]{herbst2024adrc}; aquí se ve que, además, reduce el sobrepico que deja el retardo, por lo
que el sistema responde mucho mejor ante las discrepancias con la planta.

Ahora bien, para el retardo no sirve aumentar la relación $\omega_o/\omega_c$ subiendo $\omega_o$. El
margen de fase sube un poco con esa relación (Figura~\ref{fig:adrc_margen_fase_kESO}(a) y (c)), pero
el cruce de ganancia se mueve a frecuencias más altas, por lo que el retardo máximo $\theta_{\max}$ que soporta el lazo
baja (Tablas~\ref{tab:metricas_robustez} y~\ref{tab:adrc_frecuencia_kESO}). Es decir, si el retardo empieza a
afectar el sistema, subir $\omega_o$ para tener más margen no es lo ideal, pero bajar $\omega_c$, con
$\omega_o$ fijo, sí lo aumenta. La otra opción, si el retardo se conoce, es sincronizar el observador
retrasando la orden que le entra, como proponen Herbst y Madoński \cite[sec.~9.4]{herbst2024adrc},
\cite{herbst2013practical}: con el retardo de diseño el pico de oxígeno baja y la señal de control deja
de oscilar, a cambio de una respuesta algo más lenta (Sección~\ref{sec:impl_retardo}).

\paragraph{Manejo de la saturación.} Los esquemas de manejo de la saturación, el del ESO, que
realimenta al observador la orden aplicada \cite[sec.~9.1.1]{herbst2024adrc}, y el del GPI, con el
polinomio $E_1$ del \emph{Handbook} \cite[sec.~19.1.3]{levine2011handbook}, muestran muy buenos
resultados (Sección~\ref{sec:impl_saturacion}): el lazo no se desestabiliza, incluso ante variaciones
de los parámetros de la planta (Figura~\ref{fig:sat_adrc_gpi_variacion}), y cuando el soplador sale del
límite no se observa sobrepico y el oxígeno se mantiene en la referencia. En caso contrario, sin manejo
de la saturación, el controlador se demora más en salir del límite y muestra sobrepicos más grandes,
como el GPI sin \emph{anti-windup} de la Figura~\ref{fig:sat_arranque}. Lo que ningún esquema evita es
que, mientras el soplador está en el límite, el oxígeno suba más despacio que en el diseño
(Figura~\ref{fig:adrc_lazo_cerrado_saturacion}): el \emph{anti-windup} evita que el controlador acumule
error, pero no le da al soplador más capacidad. Con la perturbación sintética se ve lo mismo: el
soplador se satura de las 4 a las 6~h y, cuando la perturbación cede, el sistema vuelve a la consigna
sin problemas (Figura~\ref{fig:pid_adrc_perturbacion}).

\paragraph{Manejo del ruido.} Para estos controladores con enfoque de rechazo de perturbaciones se
tiene un manejo del ruido muy adecuado, ya que la función de transferencia del ADRC basado en ESO es
estrictamente propia \eqref{eq:gpi_pi_filtro}, y las funciones estrictamente propias atenúan el ruido
de alta frecuencia \cite[pp.~31--32]{chen1993analog}. Herbst y Madoński llegan a lo mismo: el
controlador del ADRC es estrictamente propio y el máximo de $\omega_o$ lo pone el ruido admisible en
$u$ \cite[ecs.~4.20 y 6.30]{herbst2024adrc}. Como se mostró en la Sección~\ref{sec:impl_ruido}, con
\eqref{eq:ruido_ft} se pueden ajustar las frecuencias que se desean filtrar y las que se desean estimar
a partir del ancho de banda del observador, y en la Figura~\ref{fig:adrc_gang_kESO}(h) se ve que, entre
más grande la relación $\omega_o/\omega_c$, más ruido dejamos pasar a $u$
(Tabla~\ref{tab:adrc_frecuencia_kESO}). Por lo que el ruido, junto con el retardo, es quien realmente
pone la cota superior al ancho de banda del observador, que es el compromiso entre atenuación del ruido
de la sonda y recuperación ante la carga que plantean Cao et al.\ \cite{cao2025cascaded}. Con la
sintonía de diseño, $\omega_o$ queda por encima de la frecuencia bajo la cual está el 99~\% de la
energía de $\bar{\xi}$ (Sección~\ref{sec:ident_espectro}), así que el observador estima casi toda la
perturbación sin abrirse más al ruido (Figura~\ref{fig:pid_adrc_perturbacion}c).

\paragraph{Discretización.} Con $T_s = 50$~ms cualquiera de los métodos (ZOH, Tustin, Euler hacia
adelante o hacia atrás) sirve (Figura~\ref{fig:disc_metodos_bode}), como anticipa Chen cuando el
periodo es suficientemente pequeño \cite[sec.~13.3.1]{chen1993analog}. Con $10$ veces el ancho de banda
del lazo cerrado, el mínimo que menciona Chen, los métodos ya se separan en la fase del cruce: Tustin
es el más cercano al continuo y ZOH el más lejano (Figura~\ref{fig:disc_metodos_bode_40}). Por eso
$10$ veces es un mínimo y conviene acercarse a las $20$ que recomienda
(Figura~\ref{fig:adrc_discreto_lazo_cerrado}). En el ADRC con ESO, además, la sintonía por anchos de
banda pasa directamente al discreto: basta con llevar los polos de $\omega_c$ y $\omega_o$ al plano $z$
\eqref{eq:disc_ganancias_wc_wo}, sin rediseñar el controlador.

\paragraph{Rechazo de perturbaciones.} Viendo los resultados de la Figura~\ref{fig:gpi_sensibilidad},
correspondiente al Bode de la sensibilidad en lazo cerrado del GPI, vemos que podemos ajustar el
rechazo con cada ganancia (Tabla~\ref{tab:gpi_frecuencia}). A baja frecuencia,
$S_{\xi}(s) \approx (k_2/k_0)\,s$ \eqref{eq:gpi_caso_S_asintotas}, por lo que la parte lenta de la
perturbación la deciden $k_2$ y $k_0$: subir $k_0$ la rechaza mejor, pero baja el margen de fase. El
pico de $|S_{\xi}|$ lo baja sobre todo $k_1$, pero a costa del retardo máximo que soporta el lazo. Es
decir, lo que se gana en rechazo se paga en robustez, igual que con $\omega_o$ en el ADRC. Con la
sintonía nominal, el pico de $|S_{\xi}|$ queda por encima de donde está la mayor parte de la energía de
la perturbación (Tabla~\ref{tab:ident_resumen}), y una perturbación de entrada de $-15~\%$ apenas mueve
el oxígeno (Figura~\ref{fig:pref_escalon}). Con la perturbación sintética el lazo también la rechaza:
el oxígeno se mantiene en la consigna, salvo en el tramo en que el soplador se satura
(Figura~\ref{fig:pid_adrc_perturbacion}).

\paragraph{Implementaciones.} En la implementación se obtuvieron buenos resultados con el controlador
de rechazo activo de perturbaciones en su forma de observador de estado y en su forma GPI, para una
planta $n = 1$ y una perturbación $m = 1$, por lo que está diseñado para seguir y rechazar señales tipo
paso con error de seguimiento cero. El lenguaje que se utilizó fue texto estructurado, uno de los más
comunes dentro de la norma IEC~61131-3 \cite{iec61131}, porque de todos es el más cercano a C. Ahora
bien, los esquemas de las Figuras~\ref{fig:adrc_discreto_retardos} y \ref{fig:disc_formas_directas}
solo tienen sumas, ganancias y retardos, que existen en todos los lenguajes de la norma
(Sección~\ref{subsec:plc_lenguajes}): para el ESO son suficientes dos retardos y tres sumas, y para el
GPI discreto seis retardos y un sumador. Por lo que se puede implementar en cualquier PLC de la norma,
con una condición: trabajar en doble precisión (LREAL), porque con $T_s = 50$~ms los polos discretos
quedan muy cerca de $z = 1$ y los coeficientes necesitan muchos decimales
(Sección~\ref{subsec:disc_gpi}). Además, al estar formulado sobre el error, el bloque recibe las mismas
señales que el PID y lo puede sustituir sin recablear el lazo, que es la motivación del ADRC basado en
el error de Madoński et al.\ \cite{madonski2019general}.

\paragraph{ADRC con ESO y GPI.} Al comparar los controladores con enfoque de rechazo activo de
perturbaciones, ya sea utilizando un ESO o el GPI, se ve que tienen una relación directa. En el caso
del observador de estado extendido se busca estimar la perturbación y el error de seguimiento mediante
el observador, y el GPI trata de realizar la misma idea bajo reconstrucciones de la señal, a partir de
integraciones de la señal de control y del error \eqref{eq:gpi_integrales}. Todo esto indica que son
estructuras equivalentes que tratan de realizar la misma tarea bajo enfoques diferentes, en la línea de
las equivalencias que recopilan Rojas y Cortés-Romero entre el observador GPI, el ESO y el observador
de perturbación clásico \cite{rojas2020equivalence}. Con las ganancias $k_0$, $k_1$ y $k_2$ queda más
fácil interpretar cómo es la sensibilidad del rechazo de perturbaciones en lazo cerrado
(Figura~\ref{fig:gpi_sensibilidad}) y cómo cada ganancia afecta el sobrepico o el tiempo de
asentamiento (Tabla~\ref{tab:gpi_transitorio}). Mientras que el ADRC con ESO ayuda a interpretar el
rechazo activo de perturbaciones por medio de anchos de banda, donde se ve qué polos están dedicados al
rechazo y cuáles al control y a la respuesta transitoria \eqref{eq:sintonia_polos}.

Además, el polinomio de la sensibilidad en lazo cerrado, $e_y/\xi$, corresponde a los polos del
controlador y del observador juntos \eqref{eq:gpi_producto_S}. Por lo que vamos a tener la misma
sensibilidad, las mismas funciones de transferencia (Figura~\ref{fig:adrc_gang_kESO}), los mismos
márgenes de fase y de ganancia y los mismos anchos de banda, y lo que realmente cambia entre los
esquemas es la forma de sintonizar y su esquema de saturación. Pero incluso ahí se puede pasar de un
esquema a otro y tener el mismo resultado: el \emph{Handbook} recomienda construir $E_1$ con los polos
más rápidos del lazo cerrado \cite[ec.~19.3]{levine2011handbook}, y al tomar
$E_1(s) = (s + \omega_o)^2$ \eqref{eq:aw_gpi_wo}, los polos del observador, el GPI se comporta igual que
el ADRC también cuando satura (Figuras~\ref{fig:sat_arranque} y \ref{fig:gpi_discreto_vs_adrc}).

\paragraph{Comparación con el PID}
\label{sec:por_que}

Ahora bien, el ADRC también se puede escribir de una forma relacionada con el PID. En la literatura un
ADRC lineal se reescribe como un controlador industrial estándar \cite{stankovic2024pid} o como un PID
precedido por filtros en cascada \cite{cao2025cascaded}, así como en este caso un PI con filtro del
error \eqref{eq:gpi_pi_filtro}. Pero su diferencia con el PID ISA con filtro en la derivada radica en
que la función de transferencia de ese PID es bipropia, y la del ADRC estrictamente propia. Esa es una
primera diferencia, donde el PID tiene peores resultados en su manejo del ruido: como vimos en la
Sección~\ref{sec:impl_ruido}, al ser bipropio deja pasar el ruido hacia la señal de control, con una
$\sigma_u$ unas veinte veces mayor que la del ADRC (Tabla~\ref{tab:pid_escalon}), haciendo mayor
desgaste en los actuadores, a pesar de que la misma planta se comporte como un filtro.

Ahora bien, el PID se sintonizó para que su señal de control no salte más que la del ADRC y con la
misma sensibilidad máxima sin retardo, y aun así queda por detrás: es más lento, se pasa más en los
cambios de consigna y rechaza peor las perturbaciones (Tablas~\ref{tab:pid_escalon}
y \ref{tab:pid_pert_escalon}). En el PID una sola ganancia fija el rechazo, el esfuerzo inicial y el
ruido, mientras que en el ADRC el rechazo lo pone el observador. Sin embargo, para una planta de primer
orden como esta, un PI sintonizado con el mismo procedimiento y el mismo margen de fase del ADRC lo
supera en rechazo (Figura~\ref{fig:pid_mejores}). Lo que el ADRC sigue ganando es una señal de control que arranca suave y casi
sin ruido, y la posibilidad de subir el grado del modelo de la perturbación sin cambiar su estructura.
Además, el seguimiento y el rechazo se entienden y se sintonizan por separado, con dos anchos de banda.

Otra diferencia está en la sintonía. Un PID también admite modificar la dinámica moviendo los polos
teniendo un modelo de la planta, pero sus ganancias $K_p$, $T_i$ y $T_d$ no se relacionan directamente
con la dinámica deseada ni con la perturbación. Además, al ser de parámetros fijos no se adapta a las
variaciones de carga y de caudal \cite{qiao2016self}: queda sintonizado para un punto del rango que
recorre la planta, y en los arranques medidos del PID de la PTAR unas veces llega limpio a la consigna
y otras oscila (Sección~\ref{sec:ident_pid}). El ADRC, en cambio, se puede sintonizar por anchos de
banda \cite{gao2003scaling}, por \emph{pole placement} con las ganancias $k_0$, $k_1$ y $k_2$ para
ajustar la sensibilidad en lazo cerrado, o en su forma clásica de PI más filtro.

\section{Limitaciones del estudio}
\label{sec:limitaciones}

\begin{enumerate}
  \item \textbf{No hay trazas de la planta operando con ADRC.} Dados los tiempos y las restricciones,
  es complicado tener una ventana de tiempo para realizar la implementación del código en campo dentro
  de las fechas admisibles para el proyecto. Por eso la validación es en simulación, con la planta
  identificada, y se deja todo el proyecto probado y verificado para su puesta en marcha cuando se tenga
  el espacio en campo para instalarlo. Por la misma razón no se midió el consumo eléctrico de los
  sopladores: $E_u$ (Tabla~\ref{tab:metricas}) es la energía de la señal de control, no un consumo en
  kWh, así que el posible ahorro energético queda por confirmar en campo.

  \item \textbf{La señal de control del PLC no quedó registrada.} La consigna que el PLC envía a los
  variadores no quedó en el histórico del SCADA, así que en M2 y M3 la identificación usa la potencia en
  porcentaje de plena carga como si fuera la señal de control (Sección~\ref{sec:ident}). Por eso los
  $b_0$ identificados están en unidades de potencia y no en las del bloque del controlador, y de ahí la
  recomendación de medirlo directamente en el PLC (Sección~\ref{sec:implicaciones}).

  \item \textbf{El retardo está acotado, no medido con precisión.} El retardo se encuentra acotado y
  obtenido visualmente a partir de una de las pruebas de uno de los tres módulos dispuestos en la
  planta, por lo que dependiendo del módulo puede variar. Lo positivo de esto es que se evaluó el
  retardo con diferentes variaciones y el controlador logró responder adecuadamente, incluso mejor con
  los acondicionamientos que se hicieron.

  \item \textbf{Las variaciones se probaron una a la vez.} En las
  Figuras~\ref{fig:adrc_variacion} y \ref{fig:sat_adrc_gpi_variacion} se cambia un parámetro de la
  planta a la vez, por lo que no todas las combinaciones quedan cubiertas, y solo se validan las
  más importantes, como $b_0$ y $\theta$. Pueden aparecer casos en los que, al variar dos
  parámetros a la vez, se tengan más restricciones. El $b_0$ más alto del rango identificado es
  $1{,}5$ veces el de diseño (Tabla~\ref{tab:ident_arx}), por debajo del margen de ganancia de
  $2{,}12$; pero si se combinan la ganancia alta y la constante de tiempo corta de M3
  ($k_{dc} = 0{,}109$~mg/(L$\cdot$\%) y $\tau = 15$~min), un caso extremo, se tendría
  $b_{0,\mathrm{planta}}/b_0 = 2{,}9$ y, con el retardo de diseño, el lazo quedaría
  inestable. En ese caso habría que reajustar las ganancias para aumentar el margen de fase y que el
  lazo soporte más retardo, además de suavizar la referencia con el prefiltro o con Bézier para
  disminuir el sobrepico ante el escalón y, si se puede, sincronizar el observador.

  \item \textbf{Limitación del orden de la perturbación.} Se escogió para el diseño una perturbación de
  orden $m = 1$, un solo integrador, por lo que al cambiar la referencia por rampas u otros tipos se
  tendrán errores de seguimiento, relativamente pequeños en función de las ganancias de los
  controladores, pero de igual forma existirá el error de seguimiento. Para estos entornos industriales
  donde se usan PID no es tanto problema, pero si llegado el caso se necesita seguir otro tipo de
  referencias de forma exacta, $e_y(\infty) = 0$, deberá modificarse cada uno de los bloques de los
  controladores implementados y diseñados en este documento.
\end{enumerate}

\section{Implicaciones prácticas}
\label{sec:implicaciones}

Para quien tenga que decidir sobre un lazo de oxígeno en servicio, los
resultados se traducen en cinco recomendaciones concretas:

\begin{enumerate}
  \item \textbf{Medir $\bo$ en las unidades del bloque, no en unidades
  físicas.} Un solo escalón registrado a la velocidad del PLC basta, y evita
  toda la cadena de conversiones --- consigna a potencia, potencia a caudal de
  aire --- que es donde se acumula el error. Además, una buena metodología es que la ganancia se obtenga
  siempre cerca del punto de operación, y tratar de llegar a él posteriormente con
  algún método que acondicione la referencia.

  \item \textbf{Dejar $\bo$ por encima del de la planta si hay duda.} Con $b_0$ sobreestimado la
  respuesta es más lenta, pero se conserva la robustez. Con $b_0$ subestimado el controlador es más
  agresivo y con retardo el lazo se acerca al límite de inestabilidad (Figuras~\ref{fig:adrc_variacion}(d)
  y~\ref{fig:adrc_margenes_retardo}).

  \item \textbf{Elegir $\wo$ con el espectro de la perturbación y el ruido de la sonda.} Se elige como
  un factor de $\wc$, entre $3$ y $10$ veces como sugieren Herbst y Madoński
  \cite[sec.~3.3]{herbst2024adrc}, por encima de donde está la energía de la perturbación, pero sin
  dejar pasar el ruido de la sonda, o no lo filtraremos y tendremos una señal de control ruidosa. Cuando
  se tengan problemas de ruido y no se quiera dañar tanto el rechazo, es más fácil cambiar el
  instrumento o hacer mejores acondicionamientos en campo que realizar un sobreajuste en el observador.

  \item \textbf{Manejar la saturación y acondicionar la referencia.} Al observador le debe entrar
  $u_{\mathrm{sat}}$ y no $u$, el GPI debe tener su esquema de anti-\emph{windup}, y los cambios de
  consigna deben pasar por Bézier o por el prefiltro, que no cambian el rechazo de perturbaciones y
  bajan el pico de $u$, de $93$ a $54~\%$ con el retardo de diseño (Figura~\ref{fig:pref_escalon}).

  \item \textbf{Migrar progresivamente.} Se puede arrancar el control con anchos de banda
  conservadores, no muy grandes, y dependiendo de su comportamiento en el transitorio y en el rechazo ir
  variando las ganancias, siguiendo el resumen de la Sección~\ref{sec:impl_resumen}. Es la misma ruta
  que proponen Stanković et al.: arrancar con un ADRC que se comporte como el PID instalado y subir
  después el ancho de banda del observador \cite{stankovic2024pid}. Como el bloque recibe las mismas
  señales que el PID, es una ruta de puesta en marcha que no necesita tantos acondicionamientos.
\end{enumerate}
